\documentclass[11pt, reqno]{amsart}

\usepackage[spanish]{babel}
\usepackage[LGR, T1]{fontenc}
\usepackage[utf8]{inputenc}

../LaTeX/general.tex
% \input{../graphics.tex}

\makeatletter
\def\emailaddrname{\textit{Correo electrónico}}
\def\subtitle#1{\gdef\@subtitle{#1}}
\def\@subtitle{}

% Metadata
\def\logo#1{\gdef\@logo{#1}}
\def\@logo{}
\def\institution#1{\gdef\@institution{#1}}
\def\@institution{}
\def\department#1{\gdef\@department{#1}}
\def\@department{}
\def\professor#1{\gdef\@professor{#1}}
\def\@professor{}
\def\course#1{\gdef\@course{#1}}
\def\@course{}
\def\coursecode#1{\gdef\@coursecode{#1}}
\def\@coursecode{}

\renewcommand{\maketitle}{
\begin{center}
	\small
	\renewcommand{\arraystretch}{1.2}
	\begin{tabular}{cp{.37\textwidth}p{0.44\textwidth}}
		% \hline
		\multirow{5}{*}{\includegraphics[height=2.0cm]{\@logo}}
	  & \multicolumn{2}{c}{ \makecell{{\bfseries \@institution} \\ \@department} } \\
	  % & \multicolumn{2}{c|}{{\bfseries\@institution} \\ \@department} \\
	  \cline{2-3}
	  & \textbf{Profesor:} \@professor & \textbf{Ayudante:} \authors \\
	  % \cline{2-3}
	  & \textbf{Curso:} \@course & \textbf{Sigla:} \@coursecode \\
	  % \cline{2-3}
	  & \multicolumn{2}{l}{ \textbf{Fecha:} \@date } \\
	  % \hline
	\end{tabular}
	\\[\baselineskip]
	% {}
	% \vspace{2\baselineskip}
	{\bfseries\Large\@title}
	\ifx\@subtitle\@empty\else
		\\[1ex]
		\large\mdseries\@subtitle
	\fi
\end{center}
}
\makeatother

\usepackage{multirow, makecell}

\usepackage[
	reversemp,
	letterpaper,
	% marginpar=2cm,
	% marginsep=1pt,
	margin=2.3cm
]{geometry}
\usepackage{fontawesome}
% \makeatletter
% \@reversemargintrue
% \makeatother

% Símbolos al margen, necesitan doble compilación
\newcommand{\hard}{\marginnote{\faFire}}
\newcommand{\hhard}{\marginnote{\faFire\faFire}}

% Dependencias para los teoremas
\usepackage{xifthen}
\def\@thmdep{}
\newcommand{\thmdep}[1]{
	\ifthenelse{\isempty{#1}}
	{\def\@thmdep{}}
	{\def\@thmdep{ (#1)}}
}
\newcommand{\thmstyle}{\color{thm}\sffamily\bfseries}

% ===== Estilos de Teoremas ==========
\newtheoremstyle{axiomstyle}
	{0.3cm}
	{0.3cm}
	{\normalfont}
	{0.5cm}
	{\bfseries\scshape}
	{:}
	{4pt}
	{\thmname{#1}\thmnote{ #3}\thmnumber{ (#2)}}
\newtheoremstyle{styleC}
	{0.5cm}
	{0.5cm}
	{\normalfont}
	{0.5cm}
	{\bfseries}
	{:}
	{4pt}
	{\thmname{#1\textrm{\@thmdep}}\thmnumber{ #2}\thmnote{ (#3)}}

% ====== Teoremas (sin borde) ===========
\theoremstyle{axiomstyle}
\newtheorem*{axiom}{Axioma}

% ====== Teoremas (sin borde) ==================
\theoremstyle{styleC}
\newtheorem{thm}{Teorema}[section]
\newtheorem{mydef}[thm]{Definición}
\newtheorem{prop}[thm]{Proposición}
\newtheorem{cor}[thm]{Corolario}
\newtheorem{lem}[thm]{Lema}
\newtheorem{con}[thm]{Conjetura}
\newtheorem*{sol}{Solución}
\newtheorem*{obs}{Observación}
\newtheorem*{ex}{Ejemplo}

\newtcbox{bluebox}[1][]{enhanced jigsaw, 
  sharp corners,
  frame hidden,
  nobeforeafter,
  listing only,
  #1} % comando para crear cajas de colores

\expandafter\let\expandafter\oldproof\csname\string\proof\endcsname
\let\oldendproof\endproof
\renewenvironment{proof}[1][\proofname]{%
  \oldproof[\scshape Demostración:]%
}{\oldendproof} % comando para redefinir la caja de la demostración
\newenvironment{hint}[1][\proofname]{%
  \oldproof[\scshape Pista:]%
}{\oldendproof} % comando para redefinir la caja de la demostración

% colores utilizados
\definecolor{numchap}{RGB}{249,133,29}
\definecolor{chap}{RGB}{6,129,204}
\definecolor{sec}{RGB}{204,0,0}
\definecolor{thm}{RGB}{106,176,240}
\definecolor{thmB}{RGB}{32,31,31}
\definecolor{part}{RGB}{212,66,66}

% ====== Diseño de los titulares ===============
\usepackage[explicit]{titlesec} % para personalizar el documento, la opción <<explicit>> hace que el texto de los titulares sea un objeto interactuable

\titleformat{\subsection}[runin]
	{\bfseries}
	{\textrm{\S}\thesubsection}
	{1ex}
	{#1.}

\setlist[enumerate,1]{label=\arabic*., ref=\arabic*} % Enumerate standards

\usepackage{tikz}
\usetikzlibrary{babel,cd}
% \DeclareMathOperator\Gr{Gr}

\DeclareMathOperator\pr{pr}

\title{Aún más sobre recubrimientos y usos de van Kampen}
\date{\DTMdate{2026-10-29}}

\author{José Cuevas Barrientos}
\email{josecuevasbtos@uc.cl}
\urladdr{https://josecuevas.xyz/teach/2026-2-ayud/}

\logo{../puc_negro.png}
\institution{Pontificia Universidad Católica de Chile}
\department{Facultad de Matemáticas}
\course{Topología Algebraica}
\coursecode{MAT2850}
\professor{Mauricio Bustamante}

\begin{document}

\maketitle

\section{Recubrimientos}
Sea $X$ un espacio topológico conexo.
Un recubrimiento $Y \to X$ se dice un \strong{$G$-recubrimiento} si existe una acción fiel propiamente discontinua $G \acts Y$ tal que $X = G \coquot Y$.
Un par de $G$-recubrimientos $Y \to X$ e $Y' \to X$ se dicen \emph{isomorfos} si existen isomorfismos de recubrimientos $Y \to Y'$ que además respetan la
acción del grupo $G$ (a estos se les suelen llamar \textquote{$G$-equivariantes}).

\begin{enumerate}
	\item \textbf{Correspondencia de Galois:}
		Sea $X$ un espacio conexo y slsc.
		Describa y demuestre una biyección entre homomorfismos $\pi_1(X, x) \to G$ y $G$-recubrimientos de $X$ (posiblemente disconexos) salvo isomorfismo.

	\item Pruebe que el grupo fundamental $\pi_1(X, x)$ de toda variedad topológica $X$ basada en cualquier punto $x$ es numerable.
		% \begin{hint}
		% 	Construya un cubrimiento numerable por abiertos contráctiles (¿por qué se puede?).
		% 	Note que dos caminos que pasan por los mismos abiertos en orden son homotópicos... bueno, no en realidad, pero casi
		% 	*guiño*.
		% \end{hint}
\end{enumerate}

\section{Aplicaciones de Seifert--van Kampen}
% Donde $*$ denota el \emph{producto libre} de grupos, descrito a continuación.
Recuerde que una \emph{presentación} de un grupo $G$ es una sucesión exacta 
\begin{tikzcd}[cramped, sep=small]
	F_2 \rar["\rho"] & F_1 \rar["\varepsilon"] & G \rar & 1
\end{tikzcd}
de grupos, donde $F_j$ son libres.
Si $S = \{ x_j \}_j \subseteq F_1$ y $T = \{ y_k \}_k \subseteq F_2$ denotan bases, entonces escribimos de manera compacta
\[
	G = \langle S : \rho(T) \rangle = \langle \{ x_j \}_j : \{ \rho(y_k) = 1 \}_k \rangle.
\]
Dado un diagrama 
\begin{tikzcd}[cramped, sep=small]
	H & \lar["\varphi"'] K \rar["\psi"] & L
\end{tikzcd}
de grupos, donde la presentación de $H = \langle \{ x_j^H \} : \{ r_k^H \}_k \rangle$ y análogamente para $L$ y $K$; definimos
\[
	H *_K L := \langle \{ x_j^H, x_{j'}^L \}_{j,j'} : \{ r_k^H = r_{k'}^L = 1 \}_{k, k'} \cup \{ \varphi(x_j^K) = \psi(x_j^K) \}_j \rangle.
\]
\begin{enumerate}[resume]
	\item Demuestre que el producto libre de grupos posee la propiedad universal de ser un coproducto fibrado (eng.\ \emph{push-out})
		del diagrama y, por ende, concluya que está bien definido.
		\begin{hint}
			Para esto puede usar el \textbf{teorema de von Dyck} que dice que dado un grupo presentado $G = \langle \{ x_j \}_j
			: \{ r_k = 1\}_k \rangle$, darse un homomorfismo $\varphi\colon G \to J$ es lo mismo que darse una función
			$\{ x_j \}_j \to J$ de modo que cada relación $r_k$ (que es una palabra en los $x_j$) se anule.
		\end{hint}

	\item
		\begin{enumerate}
			\item Pruebe que, cuando $H, L$ son grupos no triviales, entonces $H * L$ (con $K = 1$) no es abeliano.
			\item\lookup
				Si $H, L$ son no triviales, ¿será que $H *_K L$ nunca sea abeliano?
		\end{enumerate}
\end{enumerate}

En esta sección, pretendemos emplear el siguiente resultado:
\begin{thm}[Seifert--van Kampen]
	Sea $X = U \cup V$ un espacio topológico cubierto por dos abiertos $U, V$ con $U \cap V \ne \emptyset$.
	Dado $x \in U \cap V$, se tiene
	\[
		\pi_1(X, x) \cong \pi_1(U, x) *_{\pi_1(U\cap V, x)} \pi_1(V, x).
	\]
\end{thm}
\begin{enumerate}[resume]
	\item\lookright
		Calcule el grupo fundamental de la \emph{rosa de $n$ pétalos} (o la \textquote{cuña de círculos}).
		
	\item Dé una presentación para el grupo fundamental de la botella de Klein.
\end{enumerate}

% \printbibliography

\end{document}
